% Options for packages loaded elsewhere
\PassOptionsToPackage{unicode}{hyperref}
\PassOptionsToPackage{hyphens}{url}
\documentclass[
  a4paper,
]{article}
\usepackage{xcolor}
\usepackage[margin=25mm]{geometry}
\usepackage{amsmath,amssymb}
\setcounter{secnumdepth}{-\maxdimen} % remove section numbering
\usepackage{iftex}
\ifPDFTeX
  \usepackage[T1]{fontenc}
  \usepackage[utf8]{inputenc}
  \usepackage{textcomp} % provide euro and other symbols
\else % if luatex or xetex
  \usepackage{unicode-math} % this also loads fontspec
  \defaultfontfeatures{Scale=MatchLowercase}
  \defaultfontfeatures[\rmfamily]{Ligatures=TeX,Scale=1}
\fi
\usepackage[]{lmodern}
\ifPDFTeX\else
  % xetex/luatex font selection
\fi
% Use upquote if available, for straight quotes in verbatim environments
\IfFileExists{upquote.sty}{\usepackage{upquote}}{}
\IfFileExists{microtype.sty}{% use microtype if available
  \usepackage[]{microtype}
  \UseMicrotypeSet[protrusion]{basicmath} % disable protrusion for tt fonts
}{}
\makeatletter
\@ifundefined{KOMAClassName}{% if non-KOMA class
  \IfFileExists{parskip.sty}{%
    \usepackage{parskip}
  }{% else
    \setlength{\parindent}{0pt}
    \setlength{\parskip}{6pt plus 2pt minus 1pt}}
}{% if KOMA class
  \KOMAoptions{parskip=half}}
\makeatother
\usepackage{color}
\usepackage{fancyvrb}
\newcommand{\VerbBar}{|}
\newcommand{\VERB}{\Verb[commandchars=\\\{\}]}
\DefineVerbatimEnvironment{Highlighting}{Verbatim}{commandchars=\\\{\}}
% Add ',fontsize=\small' for more characters per line
\newenvironment{Shaded}{}{}
\newcommand{\AlertTok}[1]{\textcolor[rgb]{1.00,0.00,0.00}{\textbf{#1}}}
\newcommand{\AnnotationTok}[1]{\textcolor[rgb]{0.38,0.63,0.69}{\textbf{\textit{#1}}}}
\newcommand{\AttributeTok}[1]{\textcolor[rgb]{0.49,0.56,0.16}{#1}}
\newcommand{\BaseNTok}[1]{\textcolor[rgb]{0.25,0.63,0.44}{#1}}
\newcommand{\BuiltInTok}[1]{\textcolor[rgb]{0.00,0.50,0.00}{#1}}
\newcommand{\CharTok}[1]{\textcolor[rgb]{0.25,0.44,0.63}{#1}}
\newcommand{\CommentTok}[1]{\textcolor[rgb]{0.38,0.63,0.69}{\textit{#1}}}
\newcommand{\CommentVarTok}[1]{\textcolor[rgb]{0.38,0.63,0.69}{\textbf{\textit{#1}}}}
\newcommand{\ConstantTok}[1]{\textcolor[rgb]{0.53,0.00,0.00}{#1}}
\newcommand{\ControlFlowTok}[1]{\textcolor[rgb]{0.00,0.44,0.13}{\textbf{#1}}}
\newcommand{\DataTypeTok}[1]{\textcolor[rgb]{0.56,0.13,0.00}{#1}}
\newcommand{\DecValTok}[1]{\textcolor[rgb]{0.25,0.63,0.44}{#1}}
\newcommand{\DocumentationTok}[1]{\textcolor[rgb]{0.73,0.13,0.13}{\textit{#1}}}
\newcommand{\ErrorTok}[1]{\textcolor[rgb]{1.00,0.00,0.00}{\textbf{#1}}}
\newcommand{\ExtensionTok}[1]{#1}
\newcommand{\FloatTok}[1]{\textcolor[rgb]{0.25,0.63,0.44}{#1}}
\newcommand{\FunctionTok}[1]{\textcolor[rgb]{0.02,0.16,0.49}{#1}}
\newcommand{\ImportTok}[1]{\textcolor[rgb]{0.00,0.50,0.00}{\textbf{#1}}}
\newcommand{\InformationTok}[1]{\textcolor[rgb]{0.38,0.63,0.69}{\textbf{\textit{#1}}}}
\newcommand{\KeywordTok}[1]{\textcolor[rgb]{0.00,0.44,0.13}{\textbf{#1}}}
\newcommand{\NormalTok}[1]{#1}
\newcommand{\OperatorTok}[1]{\textcolor[rgb]{0.40,0.40,0.40}{#1}}
\newcommand{\OtherTok}[1]{\textcolor[rgb]{0.00,0.44,0.13}{#1}}
\newcommand{\PreprocessorTok}[1]{\textcolor[rgb]{0.74,0.48,0.00}{#1}}
\newcommand{\RegionMarkerTok}[1]{#1}
\newcommand{\SpecialCharTok}[1]{\textcolor[rgb]{0.25,0.44,0.63}{#1}}
\newcommand{\SpecialStringTok}[1]{\textcolor[rgb]{0.73,0.40,0.53}{#1}}
\newcommand{\StringTok}[1]{\textcolor[rgb]{0.25,0.44,0.63}{#1}}
\newcommand{\VariableTok}[1]{\textcolor[rgb]{0.10,0.09,0.49}{#1}}
\newcommand{\VerbatimStringTok}[1]{\textcolor[rgb]{0.25,0.44,0.63}{#1}}
\newcommand{\WarningTok}[1]{\textcolor[rgb]{0.38,0.63,0.69}{\textbf{\textit{#1}}}}
\setlength{\emergencystretch}{3em} % prevent overfull lines
\providecommand{\tightlist}{%
  \setlength{\itemsep}{0pt}\setlength{\parskip}{0pt}}
\usepackage{bookmark}
\IfFileExists{xurl.sty}{\usepackage{xurl}}{} % add URL line breaks if available
\urlstyle{same}
\hypersetup{
  hidelinks,
  pdfcreator={LaTeX via pandoc}}

\author{}
\date{}

\begin{document}

\section{Anonymous Equal-Amplitude Composite-Wave Model: Basic Axiom
System
v5}\label{anonymous-equal-amplitude-composite-wave-model-basic-axiom-system-v5}

\textbf{Date:} 2026-07-21\\
\textbf{Author:} Noriaki Kihara\\
\textbf{Status:} Definition paper\\
\textbf{Version DOI:} 10.5281/zenodo.21465429 \textbf{Concept DOI:}
10.5281/zenodo.21315735

\begin{center}\rule{0.5\linewidth}{0.5pt}\end{center}

\section{1. First Principles}\label{first-principles}

\subsection{Axiom 0: Component
Anonymity}\label{axiom-0-component-anonymity}

Classification: Axiom

Basic components are not assigned individual names, privileged axes,
privileged signs, or privileged types.

The index \texttt{n} is a convenient label and is not an intrinsic name
of the component.

Forbidden:

\begin{Shaded}
\begin{Highlighting}[]
\NormalTok{making only one specific component the negative{-}sign axis}
\NormalTok{making only one specific component the time axis}
\NormalTok{making only one specific component real{-}valued}
\NormalTok{making only one specific component complex{-}valued}
\NormalTok{giving an intrinsic name to only one specific component}
\end{Highlighting}
\end{Shaded}

\begin{center}\rule{0.5\linewidth}{0.5pt}\end{center}

\subsection{Axiom 0+: Formulation
Anonymity}\label{axiom-0-formulation-anonymity}

Classification: Axiom

At the first-principle layer, theory names, formulation names, and
readout-logic names are not assigned intrinsic names.

Equations, projections, or readout rules at the first-principle layer
must not be changed on the basis of physical names, object names,
interaction names, or existing theory classifications.

Forbidden:

\begin{Shaded}
\begin{Highlighting}[]
\NormalTok{adopting a circular{-}motion equation because it is called gravity}
\NormalTok{adopting a hyperbolic equation because it is called Coulomb force}
\NormalTok{erasing signs because it is called mass}
\NormalTok{preserving signs because it is called charge}
\NormalTok{moving a component to a time axis because it is called energy}
\NormalTok{using a different closure equation because the target is a two{-}body system}
\NormalTok{using a different closure equation because the target is a three{-}body system}
\NormalTok{using a special readout rule because the target is an ABC system}
\NormalTok{choosing a right{-}hand{-}side representation to obtain a desired physical name}
\end{Highlighting}
\end{Shaded}

Allowed:

\begin{Shaded}
\begin{Highlighting}[]
\NormalTok{different readout operations applied to the same closure condition}
\NormalTok{different symmetry breaking applied to the same closure condition}
\NormalTok{different gauge stability applied to the same closure condition}
\NormalTok{different sign preservation applied to the same closure condition}
\NormalTok{different sign erasure applied to the same closure condition}
\NormalTok{different phase{-}branch selection applied to the same closure condition}
\NormalTok{different closure target sets applied to the same closure condition}
\end{Highlighting}
\end{Shaded}

Order:

\begin{Shaded}
\begin{Highlighting}[]
\NormalTok{place the same all{-}positive{-}sign zero closure}
\NormalTok{define the readout operation first}
\NormalTok{define symmetry breaking first}
\NormalTok{define sign preservation or sign erasure first}
\NormalTok{confirm gauge stability}
\NormalTok{assign working names after readout}
\end{Highlighting}
\end{Shaded}

Target set:

\[
Q(P)=\sum_n p_n^2
\]

\[
Q(A)=0
\]

\[
Q(A+B)=0
\]

\[
Q(A+B+C)=0
\]

Differences in target set are allowed.

Changing the definition of \texttt{Q} is not allowed.

\begin{center}\rule{0.5\linewidth}{0.5pt}\end{center}

\subsection{Axiom 0.5: Dimension-Wise Central Projection
Readout}\label{axiom-0.5-dimension-wise-central-projection-readout}

Classification: Auxiliary geometric axiom

No curvature correction is introduced for readout along a single axis.

Curvature-correction candidates are introduced only when two or more
independent directions form an area, volume, closed path, or transport
mismatch.

\begin{center}\rule{0.5\linewidth}{0.5pt}\end{center}

\subsection{Axiom 1: All-Positive-Sign Zero
Closure}\label{axiom-1-all-positive-sign-zero-closure}

Classification: Axiom

\[
\sum_{n=1}^{N}x_n^2=0
\]

Expanded form:

\[
x_1^2+x_2^2+\cdots+x_N^2=0
\]

Coefficient signs:

\begin{Shaded}
\begin{Highlighting}[]
\NormalTok{all terms +}
\end{Highlighting}
\end{Shaded}

Not adopted:

\[
\sum_n |x_n|^2=0
\]

Term used for closure:

\[
x_n^2
\]

Term not used for closure:

\[
x_n\bar{x}_n
\]

Applicable states:

\begin{Shaded}
\begin{Highlighting}[]
\NormalTok{stationary closed state}
\NormalTok{stable closed state}
\end{Highlighting}
\end{Shaded}

Quasi-stationary states:

\begin{Shaded}
\begin{Highlighting}[]
\NormalTok{immediately after external influence}
\NormalTok{closure recovery process}
\NormalTok{metastable state}
\end{Highlighting}
\end{Shaded}

Temporary nonzero closure residual in a quasi-stationary state:

\[
\sum_n x_n^2\ne0
\]

After stationarization:

\[
\sum_n x_n^2=0
\]

\begin{center}\rule{0.5\linewidth}{0.5pt}\end{center}

\subsection{Axiom 2: Nontrivial
Existence}\label{axiom-2-nontrivial-existence}

Classification: Axiom

\[
\exists n,\quad x_n\neq0
\]

\begin{center}\rule{0.5\linewidth}{0.5pt}\end{center}

\subsection{Axiom 3: Projection-Axis Degeneration and
Imaginary-Direction
Readout}\label{axiom-3-projection-axis-degeneration-and-imaginary-direction-readout}

Classification: Axiom

When an intrinsic coordinate system is constructed by a projection
method, the projection-axis direction is not distinguishable as a
direction for the residents of that intrinsic coordinate system.

However, if a value in the projection-axis direction remains inside the
closure condition, that value is not erased.

When the compensating quantity in the projection-axis direction is read
as \texttt{z}, it is treated inside the intrinsic coordinate system as

\[
iz
\]

This component is unobservable as a first-order direction.

When squared inside the closure expression,

\[
(iz)^2=-z^2
\]

it appears with a reversed sign.

Therefore, an unobservable compensating quantity in the projection-axis
direction appears inside the intrinsic coordinate system as a
sign-reversed compensating quantity through squared readout.

\begin{center}\rule{0.5\linewidth}{0.5pt}\end{center}

\section{2. Consequences from the First
Principles}\label{consequences-from-the-first-principles}

\subsection{Consequence 1: Insufficiency of Real Positive-Definite
Components}\label{consequence-1-insufficiency-of-real-positive-definite-components}

Classification: Derived consequence

\[
x_n\in\mathbb{R}
\Rightarrow
x_n^2\ge0
\]

\[
\sum_n x_n^2=0
\Rightarrow
x_n=0\quad\forall n
\]

\begin{center}\rule{0.5\linewidth}{0.5pt}\end{center}

\subsection{Consequence 2: Internal Generation of Sign
Reversal}\label{consequence-2-internal-generation-of-sign-reversal}

Classification: Derived consequence

\[
x_1=A,\qquad x_2=iA
\]

\[
x_1^2+x_2^2=A^2+(iA)^2=A^2-A^2=0
\]

\[
i^2=-1
\]

\begin{center}\rule{0.5\linewidth}{0.5pt}\end{center}

\subsection{Consequence 3: Necessity of Phase
Algebra}\label{consequence-3-necessity-of-phase-algebra}

Classification: Derived consequence

\[
x_n=r_n e^{i\phi_n}
\]

\[
x_n^2=r_n^2e^{i2\phi_n}
\]

Closure condition:

\[
\sum_{n=1}^{N}r_n^2e^{i2\phi_n}=0
\]

Equal-amplitude condition:

\[
r_n=r
\]

\[
\sum_{n=1}^{N}e^{i2\phi_n}=0
\]

\begin{center}\rule{0.5\linewidth}{0.5pt}\end{center}

\subsection{Consequence 4: Type
Anonymity}\label{consequence-4-type-anonymity}

Classification: Derived consequence

\[
x_n\in\mathbb{C}\quad\forall n
\]

or

\begin{Shaded}
\begin{Highlighting}[]
\NormalTok{all components are assigned the same phase{-}algebra type}
\end{Highlighting}
\end{Shaded}

\begin{center}\rule{0.5\linewidth}{0.5pt}\end{center}

\subsection{Consequence 5: Ninety-Degree Phase
Difference}\label{consequence-5-ninety-degree-phase-difference}

Classification: Derived consequence

\[
A^2+(iA)^2=0
\]

\begin{center}\rule{0.5\linewidth}{0.5pt}\end{center}

\section{3. Anonymous Equal-Amplitude Composite
Wave}\label{anonymous-equal-amplitude-composite-wave}

\subsection{Axiom 4: Anonymous Equal-Amplitude Composite
Wave}\label{axiom-4-anonymous-equal-amplitude-composite-wave}

Classification: Axiom

\[
\Psi_N=\frac{A_0}{N}\sum_{k=1}^{N}e^{i\theta_k}
\]

\[
|\psi_k|=\frac{A_0}{N}
\]

When all phases are aligned:

\[
\max|\Psi_N|=A_0
\]

\begin{center}\rule{0.5\linewidth}{0.5pt}\end{center}

\subsection{Axiom 5: Unreadability of Individual
Phase}\label{axiom-5-unreadability-of-individual-phase}

Classification: Axiom

\[
\theta_k\rightarrow\theta_k+\alpha
\]

\begin{center}\rule{0.5\linewidth}{0.5pt}\end{center}

\subsection{Axiom 6: Phase-Difference
Readout}\label{axiom-6-phase-difference-readout}

Classification: Axiom

\[
\Delta_{ij}=\theta_i-\theta_j
\]

Number of readable pair channels:

\[
\binom{N}{2}=\frac{N(N-1)}{2}
\]

\begin{center}\rule{0.5\linewidth}{0.5pt}\end{center}

\subsection{Axiom 7: System
Identification}\label{axiom-7-system-identification}

Classification: Axiom

\[
(N,A_0,\{\Delta_{ij}\})
\]

\begin{center}\rule{0.5\linewidth}{0.5pt}\end{center}

\subsection{Axiom 8: Composite Phase
Projection}\label{axiom-8-composite-phase-projection}

Classification: Axiom

\[
S_N=\sum_k e^{i\theta_k}
\]

\[
\Psi_N=\frac{A_0}{N}S_N
\]

\[
\Theta_N=\arg\Psi_N
\]

\begin{center}\rule{0.5\linewidth}{0.5pt}\end{center}

\subsection{Axiom 9: No Requirement of Inversion
Symmetry}\label{axiom-9-no-requirement-of-inversion-symmetry}

Classification: Axiom

Inversion symmetry is not required.

\begin{center}\rule{0.5\linewidth}{0.5pt}\end{center}

\section{4. External Readout}\label{external-readout}

\subsection{Axiom 10: External-Axis
Response}\label{axiom-10-external-axis-response}

Classification: Axiom candidate

\[
M_\alpha=e^{i\alpha}
\]

\[
Y(\alpha)=\frac{A_0}{N}\sum_{k=1}^{N}e^{i(\theta_k-\alpha)}
=e^{-i\alpha}\Psi_N
\]

Continuous response:

\[
S_N(\alpha)=\operatorname{Re}\left[\frac{1}{N}\sum_{k=1}^{N}e^{i(\theta_k-\alpha)}\right]
\]

Binary response:

\[
S_N(\alpha)=\operatorname{sgn}\operatorname{Re}\left[\frac{1}{N}\sum_{k=1}^{N}e^{i(\theta_k-\alpha)}\right]
\]

\begin{center}\rule{0.5\linewidth}{0.5pt}\end{center}

\subsection{Axiom 11: External Readout
Phase}\label{axiom-11-external-readout-phase}

Classification: Axiom candidate

\[
\Delta\phi_\mu=\phi_\mu^{(\Psi)}-\phi_\mu^{(M)}
\]

\[
\mu=x,y,z,t
\]

\begin{center}\rule{0.5\linewidth}{0.5pt}\end{center}

\subsection{Axiom 12: Position-Phase
Readout}\label{axiom-12-position-phase-readout}

Classification: Axiom candidate

\[
x,y,z
\]

\begin{center}\rule{0.5\linewidth}{0.5pt}\end{center}

\subsection{Axiom 13: Time-Phase
Readout}\label{axiom-13-time-phase-readout}

Classification: Axiom candidate

\[
\phi_t\equiv\omega t_{\mathrm{read}}\pmod{2\pi}
\]

\begin{center}\rule{0.5\linewidth}{0.5pt}\end{center}

\subsection{Axiom 14: Covering Phase}\label{axiom-14-covering-phase}

Classification: Axiom candidate

\begin{Shaded}
\begin{Highlighting}[]
\NormalTok{2pi{-}periodic readout}
\NormalTok{4pi{-}periodic readout}
\NormalTok{winding{-}number{-}difference readout}
\end{Highlighting}
\end{Shaded}

\begin{center}\rule{0.5\linewidth}{0.5pt}\end{center}

\section{5. Observational Phase
Discreteness}\label{observational-phase-discreteness}

\subsection{Axiom 15: Observational Phase
Discreteness}\label{axiom-15-observational-phase-discreteness}

Classification: Axiom candidate

\[
\sum_{\mathrm{cycle}}\Delta\phi=2\pi m,\qquad m\in\mathbb{Z}
\]

\[
W=\frac{1}{2\pi}\oint d\phi,\qquad W\in\mathbb{Z}
\]

Readout-cell form:

\[
\phi_{\mathrm{obs}}=m\epsilon_\phi,\qquad m\in\mathbb{Z}
\]

Phase-area candidate:

\[
Q_\phi=\sum_{i<j}\sin^2\frac{\Delta_{ij}}{2}
\]

\begin{center}\rule{0.5\linewidth}{0.5pt}\end{center}

\section{6. Waveform and Localization}\label{waveform-and-localization}

\subsection{Working Axiom A: State
Waveform}\label{working-axiom-a-state-waveform}

Classification: Working hypothesis

\[
A(\theta)=\sum_n A_n\cos(n\theta+\phi_n)
\]

\[
Z(\theta)=\sum_n r_n e^{i(n\theta+\phi_n)}
\]

\begin{center}\rule{0.5\linewidth}{0.5pt}\end{center}

\subsection{Definition A1: Normalized Equal-Amplitude Odd-Harmonic
Localized
Wave}\label{definition-a1-normalized-equal-amplitude-odd-harmonic-localized-wave}

Classification: Definition

\[
N_h=2K-1
\]

\[
S_{N_h}(u)=\frac{1}{K}\sum_{m=0}^{K-1}\cos((2m+1)u)
\]

\[
Z_{N_h}(u)=\frac{1}{K}\sum_{m=0}^{K-1}e^{i(2m+1)u}
\]

\[
\Psi_{N_h}(u)=A S_{N_h}(u)
\]

\[
\Psi_{N_h}(u)=A Z_{N_h}(u)
\]

\[
\frac{A}{K}
\]

\begin{center}\rule{0.5\linewidth}{0.5pt}\end{center}

\subsection{Theorem A1: Representative-Amplitude
Preservation}\label{theorem-a1-representative-amplitude-preservation}

Classification: Derived theorem

\[
S_{N_h}(0)=1
\]

\[
\Psi_{N_h}(0)=A
\]

\begin{center}\rule{0.5\linewidth}{0.5pt}\end{center}

\subsection{Theorem A2: Highest Odd-Harmonic Order and Localization
Width}\label{theorem-a2-highest-odd-harmonic-order-and-localization-width}

Classification: Derived theorem

\[
K=\frac{N_h+1}{2}
\]

\[
\lambda_{N_h}=\frac{\lambda_0}{N_h}
\]

\[
\Delta x_{N_h}\propto\frac{\lambda_0}{N_h}
\]

\[
N_{h,\rm req}\sim\frac{\lambda_0}{\Delta x}
\]

\[
N_{h,\rm req}^{(t)}\sim\frac{T_0}{\Delta t}
\]

\begin{center}\rule{0.5\linewidth}{0.5pt}\end{center}

\subsection{Consequence A1: Identical Construction of Spatial and
Temporal
Localization}\label{consequence-a1-identical-construction-of-spatial-and-temporal-localization}

Classification: Derived consequence

\[
\Psi(\chi,\tau)=A S_{N_{h,\chi}}(\chi-\chi_0)S_{N_{h,\tau}}(\tau-\tau_0)e^{i\phi}
\]

\begin{center}\rule{0.5\linewidth}{0.5pt}\end{center}

\subsection{Working Axiom B:
Localization}\label{working-axiom-b-localization}

Classification: Working hypothesis

\[
S_{N_h}(u)=\frac{1}{K}\sum_{m=0}^{K-1}\cos((2m+1)u)
\]

\[
Z_{N_h}(u)=\frac{1}{K}\sum_{m=0}^{K-1}e^{i(2m+1)u}
\]

\begin{center}\rule{0.5\linewidth}{0.5pt}\end{center}

\subsection{Working Axiom C: Coupling
Entrance}\label{working-axiom-c-coupling-entrance}

Classification: Working hypothesis

The entrance to interaction is low-order base-phase alignment.

\begin{center}\rule{0.5\linewidth}{0.5pt}\end{center}

\subsection{Working Axiom D: Observation
Map}\label{working-axiom-d-observation-map}

Classification: Working hypothesis

\[
s_\alpha=\mathrm{Loc}\left[\int A(\theta)M_\alpha(\theta)d\theta\right]
\]

\begin{center}\rule{0.5\linewidth}{0.5pt}\end{center}

\subsection{Working Axiom E: Squared Correlation
Strength}\label{working-axiom-e-squared-correlation-strength}

Classification: Working hypothesis

\[
P(\alpha)=\frac{|C_\alpha|^2}{\sum_\beta |C_\beta|^2}
\]

\begin{center}\rule{0.5\linewidth}{0.5pt}\end{center}

\subsection{Working Axiom F: Shared Low-Order Base
Synchronization}\label{working-axiom-f-shared-low-order-base-synchronization}

Classification: Working hypothesis

\[
\theta_A-\theta_B=\Delta
\]

\begin{center}\rule{0.5\linewidth}{0.5pt}\end{center}

\subsection{Working Axiom G: Synchronized Demodulation
Failure}\label{working-axiom-g-synchronized-demodulation-failure}

Classification: Working hypothesis

\[
V(t)=\left|\sum_n w_n e^{i\epsilon_n(t)}\right|
\]

\[
V(t)=\left|\sum_n w_n e^{-\frac12\sigma_n^2(t)}\right|
\]

\begin{center}\rule{0.5\linewidth}{0.5pt}\end{center}

\subsection{Working Axiom H: Hierarchy}\label{working-axiom-h-hierarchy}

Classification: Working hypothesis

\begin{Shaded}
\begin{Highlighting}[]
\NormalTok{low{-}order base mode}
\NormalTok{+ high{-}order localized mode}
\end{Highlighting}
\end{Shaded}

\begin{center}\rule{0.5\linewidth}{0.5pt}\end{center}

\section{7. Readout, Memory, and Complex
Representation}\label{readout-memory-and-complex-representation}

\subsection{I/Q Readout}\label{iq-readout}

Classification: Theoretical observation

\[
z=a+ib
\]

\[
a=r\cos\theta,\qquad b=r\sin\theta
\]

\[
u=r\cos\theta+r\sin\theta
=\sqrt2 r\cos\left(\theta-\frac{\pi}{4}\right)
\]

\begin{center}\rule{0.5\linewidth}{0.5pt}\end{center}

\subsection{Harmonic-Information
Readout}\label{harmonic-information-readout}

Classification: Theoretical observation

\[
Z(\theta)=\sum_n r_n e^{i(n\theta+\phi_n)}
\]

\[
P(Z)=\operatorname{Re}(Z)+\operatorname{Im}(Z)
\]

\[
P(Z)=\sum_n \sqrt2 r_n
\cos\left(n\theta+\phi_n-\frac{\pi}{4}\right)
\]

\begin{center}\rule{0.5\linewidth}{0.5pt}\end{center}

\subsection{Preservation Side and Readout
Side}\label{preservation-side-and-readout-side}

Classification: Theoretical observation

\[
z\bar z=a^2+b^2
\]

\[
z^2=a^2-b^2+2iab
\]

\begin{center}\rule{0.5\linewidth}{0.5pt}\end{center}

\subsection{Square Registry}\label{square-registry}

Classification: Working hypothesis

\[
\mathcal{A}_0=\sum_k A_k^2
\]

\[
a_0=\sum_k |A_k|^2=\sum_k A_k\bar A_k
\]

\begin{center}\rule{0.5\linewidth}{0.5pt}\end{center}

\section{8. Statistical
Classification}\label{statistical-classification}

\subsection{Pair Phase Difference}\label{pair-phase-difference}

Classification: Definition

\[
\Delta_{ij}=\theta_i-\theta_j
\]

\[
\binom{N}{2}=\frac{N(N-1)}{2}
\]

Pair composite:

\[
\psi_i+\psi_j
=2A\cos\frac{\Delta_{ij}}{2}e^{i(\theta_i+\theta_j)/2}
\]

Pair interference strength:

\[
I_{ij}=4A^2\cos^2\frac{\Delta_{ij}}{2}
\]

\begin{center}\rule{0.5\linewidth}{0.5pt}\end{center}

\subsection{Overlap Degree}\label{overlap-degree}

Classification: Definition

\[
R_N=\left|\frac{A_0}{N}\sum_{k=1}^{N}e^{i\theta_k}\right|
\]

\[
\rho_N=\frac{R_N}{A_0}
=\frac{1}{N}\left|\sum_{k=1}^{N}e^{i\theta_k}\right|
\]

\[
0\le\rho_N\le1
\]

\begin{center}\rule{0.5\linewidth}{0.5pt}\end{center}

\subsection{Fermionic Degree}\label{fermionic-degree}

Classification: Definition

\[
F_N=
\frac{2}{N(N-1)}
\sum_{i<j}\sin^2\frac{\Delta_{ij}}{2}
\]

\[
B_N=1-F_N
\]

\[
B_N=
\frac{2}{N(N-1)}
\sum_{i<j}\cos^2\frac{\Delta_{ij}}{2}
\]

\[
F_N=\frac{N}{2(N-1)}(1-\rho_N^2)
\]

\begin{center}\rule{0.5\linewidth}{0.5pt}\end{center}

\section{9. Observation Selection and Curvature
Projection}\label{observation-selection-and-curvature-projection}

\subsection{Axiom 16: Unique Observation Selection by Complete Pairwise
Relational
Waves}\label{axiom-16-unique-observation-selection-by-complete-pairwise-relational-waves}

Classification: observation-connection axiom

This axiom is an independent observation principle not derived from
Axiom 0 or Axiom 1.

For a closed system with component count \texttt{N}, define the
unordered complete pairwise relation set without proper names as

\[
\mathcal E_N
=\left\{\{i,j\}\mid 1\le i<j\le N\right\}
\]

For each relation

\[
e=\{i,j\}\in\mathcal E_N
\]

the corresponding

\[
X_e\in\mathbb C
\]

is a physical relational wave that constitutes the system, not an
observation-device channel.

The relational wave \texttt{X\_e} is not a quantity derived from
component waves; it is an independent physical state component. The
correspondence map to component waves is not specified by this axiom.

Write the relational-wave state as

\[
X_{\mathcal E}
=\left(X_e\right)_{e\in\mathcal E_N}
\]

A permutation of individual names acts as a permutation of relational
waves and does not change the closure equation, the update rule, or the
observation selection.

Define the quadratic form of the relational-wave set as

\[
q_{\mathcal E}(X)
=\sum_{e\in\mathcal E_N}X_e^2
\]

The closure of the relational waves is written, using the compensation
quantity of Axiom 3, as

\[
\sum_{e\in\mathcal E_N}X_e^2+(iR)^2=0
\]

The transposed form

\[
q_{\mathcal E}(X)=R^2
\]

may be used.

Let the number of relational waves be

\[
M=\left|\mathcal E_N\right|
\]

and let a real-coefficient generator acting identically on all
relational waves be

\[
K\in\mathbb R^{M\times M}
\]

The relational wave \texttt{X\_e} remains complex, and \texttt{K} acts
as the same real-coefficient transformation on its real and imaginary
parts.

Write the local linear update of relational waves as

\[
\dot X=KX
\]

If this update preserves the quadratic form identically,

\[
\frac{d}{d\tau}q_{\mathcal E}(X)
=X^{\mathsf T}\left(K^{\mathsf T}+K\right)X
=0
\]

Therefore, the closure-preserving generator of relational waves
satisfies

\[
K^{\mathsf T}=-K
\]

This antisymmetry is not introduced from any particular physical name or
axis name.

A real antisymmetric generator decomposes, without additional proper
names, privileged axes, or external references, into a direct sum of
two-dimensional rotation planes and a null space.

Let the number of rotation planes be

\[
r=\frac12\operatorname{rank}K
\]

and the null-space dimension be

\[
\dim\ker K=M-2r
\]

The adoption conditions for intrinsically observable directions are as
follows.

First, rotation planes become candidates for distinct directions only
when they can be mutually distinguished by the relational quantity of
independent rotation frequencies. Distinguishability is a necessary
condition and does not by itself license adoption as spatial directions.

Second, when a normal direction is read from the plane spanned by two
linearly independent relational directions, the normal may be adopted as
a direction unique up to sign only when the normal-candidate space is
one-dimensional. For a real \texttt{d}-dimensional spatial display, the
dimension of the normal-candidate space is

\[
d-2
\]

and this condition holds only for

\[
d=3
\]

Third, the number of intrinsically observable directions is not counted
from the relational-wave count \texttt{M}, the generator rank, or the
null-space dimension. It is counted as the number of directions
satisfying the uniqueness conditions above.

The minimal example is \texttt{M=3}: when

\[
\operatorname{rank}K=2,
\qquad
\dim\ker K=1
\]

the two directions of one rotation plane and the one normal direction
uniquely determined from that plane, three directions in total, may be
adopted.

When \texttt{M=6} and three rotation planes are distinguished by
independent rotation frequencies, three directions may be adopted.

When the number of distinguishable rotation planes exceeds the range of
spatial display satisfying the second uniqueness condition, the excess
rotation planes are not adopted as spatial directions and are retained
as internal phase modes.

When multiple isomorphic rotation planes remain, or when a
normal-candidate space or null space of dimension two or more remains
and its internal directions cannot be distinguished by relational
quantities alone, no single direction among them may be selected as an
observable spatial direction. Such a residue is retained as an internal
subspace without a unique basis.

The primitive readout of each relational wave is constructed from the
phase difference of Axiom 6,

\[
\Delta_{ij}=\theta_i-\theta_j
\]

and the bilinear relation obtained from the quadratic form of Axiom 1,

\[
q(X)=\sum_n X_n^2
\]

\[
B(U,V)
=\frac{1}{2}\left(q(U+V)-q(U)-q(V)\right)
\]

Readout for three or more bodies is constructed from the set of complete
pairwise relational waves and the invariant decomposition of their
closure-preserving generator.

No independent higher-order readout rule may be added on the grounds of
physical names, object names, or dimension names.

The specific coupling strengths, phase-difference dependence, and update
period of the relational-wave generator are not specified by this axiom.

They are defined as independent relational update rules consistent with
Axiom 0, Axiom 0+, Axiom 1, and this axiom.

\begin{center}\rule{0.5\linewidth}{0.5pt}\end{center}

\subsection{Axiom 17: Curvature Reaction Projection Centered on the
Future Phase
Position}\label{axiom-17-curvature-reaction-projection-centered-on-the-future-phase-position}

Classification: dynamical projection axiom

This axiom is an independent dynamical and projection principle not
derived from Axiom 0 or Axiom 1.

Assign an order index

\[
\tau\in\mathbb{N}
\]

to the state sequence of a closed system.

In the state update

\[
X(\tau)\longrightarrow X(\tau+1)
\]

the phase position on the \texttt{tau+1} side is defined as the future
phase position of that state sequence.

When the phase progression on a rotation plane uniquely constructed
without additional names by Axiom 16 has a curvature-radius readout
\texttt{rho\_c} and an angular-velocity readout \texttt{omega}, and can
be expressed as a motion whose relational virtual rotation center is the
future phase position, the center-directed compensation that closes this
virtual rotation --- that is, the reaction to the centrifugal force ---
is mapped and read out as an observable acceleration in the
circumferential direction of progression.

The center-directed compensation is the internal representation, and the
circumferential acceleration is the external readout. This axiom defines
the projection connecting the two as a rule.

The magnitude of the acceleration is

\[
|\alpha_{\mathrm{read}}|
=\rho_c|\omega|^2
\]

Let \texttt{t\_hat\_tau} be the unit vector of the circumferential
direction of progression pointing from the current phase position toward
the future phase position. The readout acceleration is

\[
\boldsymbol{\alpha}_{\mathrm{read}}
=\rho_c|\omega|^2\,\hat{\boldsymbol{t}}_\tau
\]

Here \texttt{rho\_c}, \texttt{omega}, and \texttt{t\_hat\_tau} are
relational readout quantities without physical names in the
first-principles layer. The circumferential direction is not an absolute
axis of a background space; it is determined each time from the relation
among the current phase position, the future phase position, and the
closure radius.

This axiom may not be modified on the grounds of gravity, Coulomb force,
mass, charge, or any other physical name.

The same projection rule applies to every closure mode satisfying the
same conditions.

If the phase branch opposite to the future side, a reversal of the
acceleration direction, or sign retention or sign elimination is
adopted, it must not be selected from physical names; it must be defined
in advance as an independent phase relation, branch selection, symmetry
breaking, or readout operation.

This axiom does not identify the readout acceleration with standard
gravity, the standard Coulomb force, or any known external force.

Physical names are given only after the acceleration map, distance
exponent, sign, closure preservation, and gauge stability have been
confirmed.

\end{document}
